\section{CONCLUSION}
\label{sec:conclusion}
Across the reviewed corpus, O-ISAC is organized most clearly by the physical
role of optics and the mechanism that couples communication and sensing. Optics
may carry the user link, illuminate the target, distribute a carrier, or
preserve a coherent reference, and each role creates a different set of
constraints. Shared hardware, carriers, waveforms, resources, links, and design
rules then connect the two functions at different points in the system. Their
engineering value depends on how the task, metric definition, measurement
plane, operating condition, and validation setting are reported together.

The evidence indicates that recurring tradeoff mechanisms can be compared when
their physical context and measurement contracts align. Laboratory and
prototype studies provide a broad experimental base, while field results
illustrate how topology, geometry, propagation, calibration, and service
continuity enter real systems. Application and 6G relevance become more
informative when both functional outcomes remain traceable to the same network
role and condition set. This makes transferable joint evidence the central
measure of progress and keeps individual performance values tied to their
source conditions.

Progress therefore depends on measurement contracts tailored to each modality, paired
validation under realistic disturbances, and versioned artifacts that preserve
configuration, calibration, data, and processing. These practices provide a
path from individual demonstrations to reusable benchmarks and evaluation at
network scale. They also allow subsequent designs to build on prior evidence while
preserving the physical meaning of each optical platform.
